기룡이는 하이-테크 편의점에서 아르바이트를 하고 있다. 
기룡이의 가장 중요한 업무는 매일 쌓여 있는 삼각김밥을 유통기한이 임박한 순서대로 한 줄로 정렬하는 것이다.

기룡이는 하나의 삼각김밥을 택해 맨 앞으로 보내는 작업을 반복해 삼각김밥을 정렬한다. 
그리고 이 삼각김밥 매대엔 기룡이를 도와 정렬을 대신해 주는 하이-테크 기계팔이 하나 달려 있다. 
이 기계팔은 기룡이와는 반대로 삼각김밥 중 하나를 택해 맨 뒤로 보내는 작업을 반복하여 삼각김밥을 정렬한다. 
서로 다른 작업 방식 때문에 괜한 혼란을 빚을 수 있어, 기룡이와 기계팔이 함께 정렬하는 일은 없다.

하이-테크 기계팔답게 기계팔은 항상 가능한 한 적은 선택으로 정렬하며, 이에 실직의 위협을 느낀 기룡이는 지지 않기 위해 마찬가지 최소 횟수로 정렬할 수 있게 되었다.

그렇게 열심히 일하던 기룡이는 같은 순서로 배열된 삼각김밥이라도 자신과 기계팔이 이를 정렬하는 선택의 횟수가 서로 다를 수 있다는 것을 깨달았다. 
더 나아가 기룡이는 $N$개의 삼각김밥에 대해 최선을 다했을 때, 기룡이는 정확히 $A$번, 기계팔은 정확히 $B$번 삼각김밥을 택해 정렬할 수 있는 초기 삼각김밥의 배치가 몇 가지나 될지 궁금해졌다. 

하지만 기룡이는 지금 삼각김밥을 정리하느라 바쁘다! 기룡이를 도와 이를 구하는 프로그램을 작성해 보자.